\subsection{Edge-IQA Validation}
\label{sec:tcm-iqa}

On the validation split ($n{=}572$ clear + synthetic blur pairs), Fig.~\ref{fig:iqa-hist} shows overlapping $Q_{\mathrm{img}}$ histograms.
Table~\ref{tab:iqa-calibration} lists cohort medians ($0.516/0.414$), deployed $\tau_{B2}{=}0.465$, and negative separability ($-0.541$): some blurred validation crops score above the weakest clear crops, so $\tau$ is a cohort-median operating point rather than a hard separator.
Spearman $\rho{\approx}0.36$--$0.47$ on real texture motivates pairing interpretable B2 with learned comparators (Table~\ref{tab:learned-hybrid}).

% Table: ShezhenV3-COCO — auto-generated
\begin{table}[t]
\centering
\caption{ShezhenV3-COCO tongue dataset splits used in Paper~I.}
\label{tab:dataset}
\small
\resizebox{\columnwidth}{!}{%
\begin{tabular}{lrr}
\hline
Split & Images & Purpose \\
\hline
Train & 5594 & held out (future fine-tuning) \\
Validation & 572 & Edge-IQA $\tau$ calibration \\
Test & 553 & main matrix (B0--B3) \\
\hline
Total & 6719 & COCO-format bounding boxes \\
\hline
\end{tabular}%
}
\end{table}


\begin{table}[t]
\centering
\caption{Edge-IQA calibration on ShezhenV3 validation split ($n=572$ clear + $n=572$ synthetic blur).}
\label{tab:iqa-calibration}
\small
\resizebox{\columnwidth}{!}{%
\begin{tabular}{lcc}
\hline
Metric & Clear cohort & Blur cohort \\
\hline
Median $Q_{\mathrm{img}}$ & 0.516 & 0.414 \\
Mean $Q_{\mathrm{img}}$ & 0.519 & 0.414 \\
Std & 0.116 & 0.102 \\
\hline
Routing threshold $\tau$ & \multicolumn{2}{c}{0.465} \\
Spearman $\rho$ (M5) & \multicolumn{2}{c}{0.335 ($n=1144$)} \\
\hline
\end{tabular}%
}
\end{table}


Each test frame is scored by the same \texttt{scorer.py} invoked offline; latency components follow \S\ref{sec:setup}.
This ensures Table~\ref{tab:main-results} reflects \emph{real} quality scores while preserving the reproducible closed-loop simulator for RTT aggregation.
An auxiliary Franka track (M6, 50 trials on fake hardware) logs JSONL through the live gateway stack and confirms routing decisions are reproducible on the execution path; M6 RTT trends are reported in the online supplement and must not be merged into Table~\ref{tab:main-results}.

\begin{figure*}[!t]
  \centering
  \includegraphics[width=0.92\linewidth]{fig_iqa_hist.pdf}
  \caption{Edge-IQA $Q_{\mathrm{img}}$ on ShezhenV3 validation: clear vs.\ synthetically blurred pairs ($n{=}1{,}144$). Overlap motivates cohort-median $\tau$ rather than hard separation.}
  \label{fig:iqa-hist}
\end{figure*}
